\subsection{布饶尔群的应用}

在本小节中，K 是一个体，L 是 K 上的一个伽罗瓦代数，带作用 \(\lambda: G \to \operatorname{Aut}_{K}(L)\)。我们用 n 表示 L 在 K 上的次数。我们有 \(n = \operatorname{Card}(G)\)。

定理 3. —— 设 \(\mathcal{E} = (\Gamma, \iota, \pi)\) 是 G 被 L* 的一个 \(\tau\)-扩张。代数 \(\mathbf{A}[\mathcal{E}; \mathrm{L}]\) 是中心单的，且在 K 上的次数为 \(n^2\)。此外，典范同态 \(u\)（从 L 到 \(\mathbf{A}[\mathcal{E}; \mathrm{L}]\) 的）是从 L 到 \(\mathbf{A}[\mathcal{E}; \mathrm{L}]\) 的一个极大交换子代数上的同构。

引理 12. —— a) 除 \(\{0\}\) 与 L 外，不存在 L 的任何在 G 下不变的理想。

\begin{enumerate}
\def\labelenumi{\alph{enumi})}
\setcounter{enumi}{1}
\tightlist
\item
  设 g 是 G 的一个异于单位元素的元素，并设 \(\mathfrak{a}_{g}\) 是 L 的由形如 \(x - \lambda(g)x\) 的元素（其中 x 取遍 L）所生成的理想。我们有 \(\mathfrak{a}_{g} = L\)。
\end{enumerate}

设 \(\mathscr{S}\) 是 L 的极大理想所成之集；对 \(\mathscr{S}\) 的任一子集 S，令 \(\mathfrak{a}(S) = \bigcap_{\mathfrak{m} \in S} \mathfrak{m}\)。由于环 L 是交换且半单的（VIII, 第 310 页，注 6），映射 \(W \mapsto \mathfrak{a}(W)\) 是从 \(\mathscr{S}\) 的子集所成之集到 L 的理想所成之集上的一个双射（VIII, 第 142 页，命题 9）。理想 \(\mathfrak{a}(S)\) 在 G 下不变当且仅当 S 在 G 下不变。由于 G 传递地作用于 \(\mathscr{S}\)（VIII, 第 308 页，定理 2, vi)），\(\mathscr{S}\) 的在 G 下不变的子集只有 \(\varnothing\) 与 \(\mathscr{S}\)。由于 \(\mathfrak{a}(\varnothing) = L\) 且 \(\mathfrak{a}(\mathscr{S}) = \{0\}\)，我们便证明了断言 a)。

我们用反证法证明 b)；假设我们有 \(\mathfrak{a}_g\neq \mathrm{L}\)。设 \(\mathfrak{m}\) 是 L 的一个包含 \(\mathfrak{a}_g\) 的极大理想。于是对 L 中的每个 x，我们有 \(\lambda (g)x\equiv x\bmod \mathfrak{m}\)。特别地，m 在 g 下不变，且 \(\lambda(g)\) 在体 L/m 上由恒同映射诱导。由 VIII, 第 308 页定理 2 的性质 (vi)，\(G_{m}\) 的元素 g 是平凡的，这与 b) 的假设矛盾。

我们现在证明定理 3。向量空间 \(\mathbf{A}[\mathcal{E};\mathrm{L}]\) 具有有限维数 \(\operatorname{Card}(\mathrm{G})[\mathrm{L}:\mathrm{K}] = n^2\)（在 \(\mathrm{K}\) 上；由 VIII, 第 316 页以及所设等式 \(\operatorname{Card}(\mathrm{G}) = [\mathrm{L}:\mathrm{K}] = n\)）。

设 \(\mathfrak{a}\) 是代数 \(\mathbf{A}[\mathcal{E};\mathrm{L}]\) 的一个非零双边理想。我们采用 VIII, 第 316 页的记号。每个元素 \(a\)（\(\mathbf{A}[\mathcal{E};\mathrm{L}]\) 的）都可以唯一地写成 \(a = \sum_{g\in \mathrm{G}}a_g\varepsilon_g\)，其中 \(a_g\in \mathrm{L}\) 对每个 \(g\in \mathrm{G}\) 成立；我们用 \(\Phi (a)\) 表示元素 \(g\)（\(\mathrm{G}\) 的、使得 \(a_g\neq 0\) 的）所成的集合。由 VIII, 第 316 页的公式 (32)，我们有关系式

\[
\varepsilon_ {g} \varepsilon_ {g ^ {\prime}} = c _ {\sigma} (g, g ^ {\prime}) \varepsilon_ {g g ^ {\prime}}\tag{38}
\]

对所有 \(g, g' \in \mathrm{G}\) 成立，由此推出

\[
\Phi (a \varepsilon_ {g}) = \Phi (a). g\tag{39}
\]

对每个 \(g\in \mathbf{G}\) 与每个 \(a\in \mathbf{A}[\mathcal{E};\mathrm{L}]\)。

设 \(a\) 是 \(\mathfrak{a}\) 的一个非零元素，\(\Phi(a)\) 对于包含关系为极小；由公式 (39)，如有必要，把 \(a\) 换成形如 \(a\varepsilon_{g^{-1}}\) 的元素（其中 \(g\in\Phi(a)\)），我们就可以假设我们有 \(e\in\Phi(a)\)。设 \(s\) 是 \(\Phi(a)\) 的一个异于 \(e\) 的元素。由引理 12, b)，存在 L 的一个元素 \(x\)，使得 \(a_s(x-\lambda(s)\cdot x)\neq0\)。但我们有关系式

\[
x a - a x = \sum_ {g} a _ {g} (x - \lambda (g) \cdot x) \varepsilon_ {g},\tag{40}
\]

其中求和取遍 \(\Phi(a)\) 中异于 e 的元素 g。由于 \(\Phi(a)\) 的极小性，我们得到 xa - ax = 0，但这与对 x 的选取矛盾。于是我们用反证法证明了 \(\Phi(a)\) 只含 G 的单位元素 e，从而我们有 \(a \in L\)。

于是 \(L \cap a\) 是 L 的一个不化归为 \(\{0\}\) 的理想。此外，对 L 中的每个 x，我们有 \(\varepsilon_{g}x\varepsilon_{g}^{-1} = \lambda(g) \cdot x\)，故 \(L \cap a\) 在 G 下不变。由引理 12，我们因此有 \(L \cap a = L\)，即 \(L \subset a\)。由于 L 含有 \(A[E;L]\) 的单位元素，我们有 \(a = A[E;L]\)。

由于代数 \(A[E;L]\) 在 K 上具有非零的有限维数，且它仅有的双边理想是 \(\{0\}\) 与 \(A[E;L]\)，故它是单的。

我们证明与 L 交换的每个元素 \(a\)（\(\mathbf{A}[\mathcal{E};\mathrm{L}]\) 的）都属于 L。我们把 \(a\) 写成 \(\sum_{g\in \mathrm{G}}a_g\varepsilon_g\)，其系数 \(a_{g}\) 在 L 中。对 L 中的每个 \(x\)，我们有 \(xa - ax = 0\)，而关系式 (40) 表明我们有

\[
a _ {g} (x - \lambda (g) \cdot x) = 0\tag{41}
\]

对每个 \(g \in \mathrm{G}\) 与 \(x \in \mathrm{L}\)。由引理 12，我们因此有 \(a_g = 0\)（对 \(g \neq e\)）；从而 \(a \in \mathrm{L}\)。

最后，我们来确定 \(\mathbf{A}[\mathcal{E};\mathrm{L}]\) 的中心。若 \(z\) 属于中心，则它与 \(\mathrm{L}\) 交换，从而属于 \(\mathrm{L}\)。但这时我们有

\[
0 = \varepsilon_ {g} z - z \varepsilon_ {g} = (\lambda (g) \cdot z - z) \varepsilon_ {g}
\]

对 G 的每个 g 成立。于是 z 在 L 的自同构群 \(\lambda(\mathrm{G})\) 下不变。因此按假设我们有 \(z \in K\)。

定理 4. —— 设 A 是 K 上的一个有限次数中心单代数，并设 L 是 A 的一个极大交换子代数。则存在一个 \(\tau\)-扩张 \(\mathcal{E}\)（G 被 \(\mathrm{L}^*\) 的），使得 A 同构于 \(\mathbf{A}[\mathcal{E};\mathrm{L}]\)。

不失一般性，我们可以假设 L 是 A 的一个极大交换子代数。设 \(\Gamma\) 是由 A 中满足如下条件的可逆元素 \(\gamma\) 组成的乘法群：存在 G 中的 g 使得

\[
\gamma x \gamma^ {- 1} = \lambda (g). x\tag{42}
\]

对 L 的每个元素 x 成立。若给定了 \(\gamma \in \Gamma\)，则满足这一关系的元素 g 是唯一的；我们把它记作 \(\pi(\gamma)\)。直接看出，\(\pi\) 是从 \(\Gamma\) 到 G 的一个同态，且其核等于 \(L^{*}\)。\(\pi\) 的满射性由斯科伦--诺特定理（VIII, 第 263 页，推论）得出。

如果把从 \(L^{*}\) 到 \(\Gamma\) 中的典范单射记作 \(\iota\)，那么由诸构造推出，\(\mathcal{E} = (\Gamma, \iota, \pi)\) 是一个 \(\tau\)-扩张（G 被 \(L^{*}\) 的）。设

\[
u: \mathrm{L} \to \mathbf {A} [ \mathcal {E}; \mathrm{L} ] \quad \text { and } \quad v: \Gamma \to \mathbf {A} [ \mathcal {E}; \mathrm{L} ]
\]

为诸典范同态。由 \(\mathbf{A}[\mathcal{E};\mathrm{L}]\) 的普遍性质（VIII, 第 315 页，命题 12），存在唯一的代数同态 \(f:\mathbf{A}[\mathcal{E};\mathrm{L}]\to \mathrm{A}\)，使得 \(f\circ u(x) = x\) 与 \(f\circ v(\gamma) = \gamma\)（对 \(x\in \mathrm{L}\) 与 \(\gamma \in \Gamma\)）。由于代数 \(\mathbf{A}[\mathcal{E};\mathrm{L}]\) 是单的，同态 \(f\) 是单射的。而代数 \(\mathbf{A}[\mathcal{E};\mathrm{L}]\) 的次数为 \(n^2\)（在 \(\mathbf{K}\) 上），\(\mathrm{A}\) 亦然，因为 \(\mathrm{L}\) 是 \(\mathrm{A}\) 的一个极大交换半单子代数，且我们有 \(n = [\mathrm{L}:\mathrm{K}]\)（VIII, 第 262 页，命题 3, (ii)）。于是 \(f\) 是双射。

定义 2. —— 设 \(\mathscr{S}\) 是 L 的极大理想所成之集。我们定义

\[
\operatorname{Br} (\mathrm{L} / \mathrm{K}) = \bigcap_ {\mathfrak {m} \in \mathscr {S}} \operatorname{Ker} (r _ {(\mathrm{L} / \mathfrak {m}) / \mathrm{K}}),
\]

其中 \(r_{(L/\mathfrak{m})/K} : \mathrm{Br}(K) \to \mathrm{Br}(L/\mathfrak{m})\) 是标量扩张同态（VIII, 第 281 页）。

定理 5. —— 存在一个群同构

\[
\Psi : \mathrm{Ex} _ {\tau} (\mathrm{G}, \mathrm{L} ^ {*}) \longrightarrow \mathrm{Br} (\mathrm{L} / \mathrm{K})
\]

它把一个 \(\tau\)-扩张 \(\mathcal{E}\)（G 被 \(\mathrm{L}^*\) 的）的类映为 \(\mathrm{Br}(\mathrm{L} / \mathrm{K})\) 中代数 \(\mathbf{A}[\mathcal{E};\mathrm{L}]\) 的类。

为了定义 \(\Psi\) 并验证它是双射，我们必须确立以下几点：

\begin{enumerate}
\def\labelenumi{\alph{enumi})}
\item
  若 \(\mathcal{E}\) 与 \(\mathcal{E}'\) 是两个同构的 \(\tau\)-扩张（\(G\) 被 \(L^*\) 的），则代数 \(\mathbf{A}[\mathcal{E}; L]\) 与 \(\mathbf{A}[\mathcal{E}'; L]\) 同构。
\item
  反之，若代数 \(\mathbf{A}[\mathcal{E};\mathrm{L}]\) 与 \(\mathbf{A}[\mathcal{E}';\mathrm{L}]\) 同构，则 \(\tau\)-扩张 \(\mathcal{E}\) 与 \(\mathcal{E}'\)（\(\mathrm{G}\) 被 \(\mathrm{L}^*\) 的）同构。
\item
  在 \(\mathrm{Br(L / K)}\) 的每个类中，都有一个代数 \(E\)，它包含 \(L\) 作为极大交换子代数。
\item
  若 \(\mathrm{E}\) 是 \(\mathrm{K}\) 上的一个有限次数中心单代数，它包含 \(\mathrm{L}\) 作为极大交换子代数，则存在一个 \(\tau\)-扩张 \(\mathcal{E}\)（\(\mathrm{G}\) 被 \(\mathrm{L}^*\) 的），使得 \(\mathrm{E}\) 同构于 \(\mathbf{A}[\mathcal{E};\mathrm{L}]\)。
\end{enumerate}

断言 a) 由交叉乘积的构造得出；b) 由 VIII, 第 263 页的推论得出。断言 c) 由 VIII, 第 281 页的命题 5 得出，而 d) 就是上面的定理 4。

剩下要验证 \(\Psi\) 是一个群同态；为此，只需证明：若 \(\mathcal{E}_1 = (\Gamma_1,\iota_1,\pi_1)\) 与 \(\mathcal{E}_2 = (\Gamma_2,\iota_2,\pi_2)\) 是两个 \(\tau\)-扩张，则代数 \(\mathbf{A}[\mathcal{E}_1\mathcal{E}_2;\mathrm{L}]\) 等价于代数 \(\mathbf{A}[\mathcal{E}_1;\mathrm{L}]\otimes \mathbf{A}[\mathcal{E}_2;\mathrm{L}]\)。我们把积扩张 \(\mathcal{E}_1\mathcal{E}_2\) 记作 \(\mathcal{E} = (\Gamma ,\iota ,\pi)\)。群 \(\Gamma\) 同构于同态 \(\rho\)（从 \(\mathrm{L}^*\) 到纤维积 \(\Gamma_1\times_{\mathrm{G}}\Gamma_2\)，把 \(\mu\) 映为 \((\iota_1(\mu),\iota_2(\mu)^{-1})\)）的余核。对 \(i\in \{1,2\}\)，记 \(\mathrm{A}_i = \mathbf{A}[\mathcal{E}_i;\mathrm{L}]\)。设 \(u_{i}:\mathrm{L}\to \mathrm{A}_{i}\) 与 \(v_{i}:\Gamma_{i}\rightarrow \mathrm{A}_{i}^{*}\) 为诸典范同态。我们通过诸同态 \(u_{i}\) 把 L 与它的像等同起来，这些同态使 L 成为诸 \(\mathrm{A}_i\) 的极大交换子代数。我们用 \(\mathrm{V}_i\) 表示由 \(\mathrm{A}_i\) 中的左乘定义的 L-向量空间。环 \(\mathrm{L}\otimes_{\mathrm{K}}\mathrm{A}_i^{\circ}\) 作用于 \(\mathrm{V}_i\)。由于它是单的并且在 L 上的维数为 \(n^2\)，我们得到一个从 \(\mathrm{L}\otimes_{\mathrm{K}}\mathrm{A}_i^{\circ}\) 到 \(\operatorname {End}_{\mathrm{L}}(\mathrm{V}_{i})\) 上的同构。于是环 \(\mathrm{L}\otimes_{\mathrm{K}}\mathrm{A}_1^{\circ}\otimes_{\mathrm{K}}\mathrm{A}_2^{\circ}\)（我们可以把它与 \((\mathrm{L}\otimes_{\mathrm{K}}\mathrm{A}_1^{\circ})\otimes_{\mathrm{L}}(\mathrm{L}\otimes_{\mathrm{K}}\mathrm{A}_2^{\circ})\) 等同）同构于 \(\operatorname {End}_{\mathrm{L}}(\mathrm{V}_{1}\otimes_{\mathrm{L}}\mathrm{V}_{2})\)。令 \(\mathrm{C} = \operatorname {End}_{\mathrm{A}_1^{\circ}\otimes_{\mathrm{K}}\mathrm{A}_2^{\circ}}(\mathrm{V}_1\otimes_\mathrm{L}\mathrm{V}_2)\)。由 VIII, 第 282 页的引理 3，环 C 相似于 \(\mathrm{A}_1\otimes_{\mathrm{K}}\mathrm{A}_2\)，且 \(\mathrm{L}\otimes 1\otimes 1\) 是 C 的一个极大交换子代数。对每一对 \((\gamma_{1},\gamma_{2})\in \Gamma_{1}\times \Gamma_{2}\)（满足 \(\pi_1(\gamma_1) = \pi_2(\gamma_2)\)），我们用 \(w(\gamma_1,\gamma_2)\) 表示

唯一的 \(\lambda (\pi_1(\gamma_1))\)-半线性的（II, §1, No.~13, 第 223 页）自同态，使得

\[
w (\gamma_ {1}, \gamma_ {2}) (x _ {1} \otimes x _ {2}) = v _ {1} (\gamma_ {1}) x _ {1} \otimes v _ {2} (\gamma_ {2}) x _ {2}
\]

对 \(x_{1} \in V_{1}\) 与 \(x_{2} \in V_{2}\)。我们有 \(w(\gamma_{1}, \gamma_{2}) \in \mathrm{C}^{*}\)，且 w 是从纤维积 \(\Gamma_{1} \times_{G} \Gamma_{2}\) 到 \(C^{*}\) 的一个群同态。这个同态在 \(\rho\) 的像上是平凡的，并诱导一个从 \(\Gamma\) 到 \(C^{*}\) 的同态 v。用 \(u : L \to C\) 表示由 \(u : \ell \mapsto \ell \otimes 1 \otimes 1\) 给出的态射。我们可以验证关系式

\[
u (\alpha) = v (\iota (\alpha)) \quad \text { and } \quad u (^ {\gamma} \beta) = v (\gamma) u (\beta) v (\gamma) ^ {- 1}
\]

对 \(\alpha \in \mathrm{L}^*\)、\(\beta \in \mathrm{L}\) 与 \(\gamma \in \Gamma\)。VIII, 第 315 页的命题 12 给出一个同态 \(f\)（从代数 \(\mathbf{A}[\mathcal{E};\mathrm{L}]\) 到 \(\mathrm{C}\) 的）。由于代数 \(\mathbf{A}[\mathcal{E};\mathrm{L}]\) 是单的，同态 \(f\) 是单射的。代数 \(\mathrm{C}\) 与 \(\mathbf{A}[\mathcal{E};\mathrm{L}]\) 在 \(\mathrm{K}\) 上有相同的维数，故 \(f\) 是同构。

注. —— 1) 若 L 是 K 上的一个平展代数且 G 是 L 的自同构群，则代数 \(A[E;L]\) 并不总是中心单的（例如，可取 \(L = K^{n}\) 与 \(G = S_{n}\)）。

\begin{enumerate}
\def\labelenumi{\arabic{enumi})}
\setcounter{enumi}{1}
\tightlist
\item
  我们可以如下计算一个 2-闭上链 c，它与一个被带群 G 的有限伽罗瓦扩张 L 分裂的代数 A 相伴。首先，存在一个 K-同态 \(\varphi: A \to \mathbf{M}_{m}(L)\)，其中 \([A:K] = m^{2}\)。对 \(g \in G\)，用 \(\varphi^{g}\) 表示从 A 到 \(\mathbf{M}_{m}(L)\) 的由 \(a \mapsto \varphi(g^{-1}ag)\) 给出的同态。由斯科伦--诺特定理（VIII, 第 256 页，定理 3），对每个 \(g \in G\)，存在元素 \(u_{g}\)（\(\mathbf{GL}_{m}(L)\) 的），使得
\end{enumerate}

\[
\varphi^ {g} (a) = u _ {g} \varphi (a) u _ {g} ^ {- 1}
\]

对 \(a\in \mathrm{A}\)。然后我们令

\[
c (g, g ^ {\prime}) = u _ {g} u _ {g ^ {\prime}} u _ {g g ^ {\prime}} ^ {- 1}.
\]

我们也可以定义 G 被 \(L^{*}\) 的一个扩张，其作法是使用 \(\varphi\)：我们考虑群 \(\Gamma \subset \mathbf{GL}_{m}(L)\)，它由满足如下条件的 \(\gamma\) 组成：存在 \(g \in G\) 使得

\[
\varphi^ {g} (a) = \gamma \varphi (a) \gamma^ {- 1}
\]

对每个 \(a \in \mathrm{A}\)。这个扩张的类是在 \(\Psi\) 之下 \(\mathrm{A}\) 的类（在 \(\mathrm{Br}(\mathrm{L} / \mathrm{K})\) 中的）的逆像。

推论. —— 映射 \(\Phi_{L/K} = \Theta \circ \Psi^{-1}\) 定义了一个从 \(\mathrm{Br}(\mathrm{L}/\mathrm{K})\) 到 \(\mathrm{H}^{2}(\mathrm{G}, \mathrm{L}^{*})\) 的群同构。

设 \(K'\) 是 K 的一个扩张，并设 \(\varphi: K' \to L\) 是一个 K-代数态射。由 G 中满足 \(\lambda(h) \circ \varphi = \varphi\) 的元素 h 组成的集 H 是 G 的一个子群，而 \(K'\)-代数 L（赋予 \(\lambda\) 到 H 的限制）是 \(K'\) 上的一个伽罗瓦代数。

命题 14. —— 下面的图交换：

\[
\begin{array}{c} \operatorname{Br} (\mathrm{L} / \mathrm{K}) \xrightarrow {\Phi_ {\mathrm{L} / \mathrm{K}}} \operatorname{H} ^ {2} (\operatorname{G}, \operatorname{L} ^ {*}) \\ \Biggl \downarrow^ {r _ {\mathrm{K} ^ {\prime} / \mathrm{K}}} \qquad \qquad \qquad \qquad \Biggl \downarrow^ {\operatorname{Res} _ {\mathrm{H}} ^ {\mathrm{G}}} \\ \operatorname{Br} (\mathrm{L} / \mathrm{K} ^ {\prime}) \xrightarrow {\Phi_ {\mathrm{L} / \mathrm{K} ^ {\prime}}} \operatorname{H} ^ {2} (\operatorname{H}, \operatorname{L} ^ {*})  . \end{array}
\]

这由 VIII, 第 299 页的命题 7 与 VIII, 第 317 页的命题 13 得出。

